\chapterimage{special.pdf}
\chapterimagecaption{\textit{Ophelia}  (1852) by John Everett Millais}

\chapter{Special Situations}

Here is a compilation of some special extraordinary circumstances that could hinder your degree progression.

\section{Academics \& Research}

\subsection{Prelim timing for direct-entry PhD students}   \index{Prelim}     
        In this subsection, ``direct-entry PhD" refers to students who join the PhD program with no previous experience in a graduate physics program---i.e. as PhD 1. Students who enter the PhD program as PhD 2, whether via the McGill fast track, with an MSc from McGill, or an MSc from a prior institution, are not considered direct-entry. 
        
        The general advice is to take the PhD preliminary exam by the end of one's first term of PhD 3 study--e.g. for a student who entered in the fall as PhD2 via fast-track, this would be the fall of the following year; or for a student who entered in the winter as PhD2 with an MSc from another institution, this would be the winter of the following year. Fortunately, this logic extends to direct-entry students---for example, a direct-entry student who joins McGill graduate physics studies in Fall 2026 must take their prelim by the end of Fall 2028.

        However, this is not well-documented on the department website's table of hypothetical prelim timing situations---possibly because direct-entry admissions are not widespread across all parts of the department. At the time of writing, several direct-entry students who embarked upon their McGill graduate physics studies in Fall 2024 and Fall 2025 successfully confirmed with the Physics GPC that they should plan on taking their prelims by the end of Fall 2026 and Fall 2027, respectively; confirming the ``end of first term of PhD 3" rule. 
        
        Please note, however, that, as of Summer 2026, MyProgress is not configured to represent this edge case well: it shows the prelim milestone as overdue for direct-entry students at the end of their fourth non-summer term (as opposed to fifth non-summer term, the GPC-confirmed policy).

    \subsection{Failing a course}\index{Failing a course}

    
    Failing a course can be scary but we promise its not the end of the world. Your options after your first failure (i.e, any grade less than a B-- or any incomplete grade such as J,K, or L) are as follows:
        \begin{itemize}
            \item Contest the grading with the professor (not much can actually be done but always worth an attempt)
            \item Write a supplemental exam upon approval by the department and professor
            \item Retake the failed course
            \item Substitute the failed course by completing an equivalent course. Neither this nor retaking the original class will replace the original failed grade on your transcript.
        \end{itemize}
        Students who fail a second course will be automatically be withdrawn from the University. 

    \subsection{Withdrawing} \index{Withdrawing}
        Harking back to the earlier section, a student will be withdrawn from the University if they:
        \begin{itemize}
            \item Fail two courses
            \item Obtain two unsatisfactory progress tracking reports and the academic unit recommends they be withdrawn
            \item Fail one course and receive one unsatisfactory progress tracking report and is recommended to be withdrawn by their academic unit.
        \end{itemize}
        
    \subsection{Death of a supervisor}
        
        In the event of a prolonged absence of your supervisor, whether its a long sabbatical, death, retirement or if your supervisor got a new job elsewhere, talking to the department and other professors to figure out an individual course of action is the best way forward. Each case is different, so reach out to the Graduate Studies Advisor and the Graduate Studies Coordinator to outline the best way forward for your specific situation.
        
\section{Extra-Curricular}

    \subsection{Paternity/Maternity Leave} \index{Paternity Leave}\index{Maternity Leave}
        A leave of absence may be requested for parental leave from McGill by filling out a request from. For a maternity or parental leave, the eligibility period of a maximum of 52 consecutive weeks is determined based on when the child is born; if the leave is interrupted for one or two terms, the eligibility period cannot be extended. Most agencies also allow you to pause your funding for a period of time while you take your parental leave. You must inform these agencies of these planned absences.


    \subsection{Sports Injuries}\index{Sports Injuries}
        Whether you're a member of one of our awesome intramural teams or were out doing your own activity or just injured yourself in your own special way, McGill Sports clinic has got you covered. Through the McGill Sports clinic you can access an array of services and clinics that best suit your needs. The list of services the sports clinic has to offer are: 
        \begin{itemize}
            \item Sports Medicine
            \item Massage Therapy
            \item Physiotherapy
            \item Concussion assessment and treatment
            \item Osteopathy
            \item Vestibular Rehabilitation
        \end{itemize}
        The sports clinic is located on the 3rd floor of the McGill Sports Centre. McGill students get preferred rates and can typically see a doctor within a couple of days of calling to make an appointment. You can make an appointment online or by calling (514)-398-7007.

    \subsection{International Taxes} \index{Taxes!International Taxes}
        International students with citizenship in most countries are all set to file taxes in Canada alone (see 
        \hyperref[tax]{\textcolor{purble}{\ref*{tax}}}). Citizens of China, Eritrea, and the United States should take note that they are responsible for filing taxes to the relevant agencies in these countries during each fiscal year where they remain a citizen---even if they did not set foot in, reside in, or earn money in said country during said fiscal year.

    \subsection{NRC Reliability Status}
        Some students will need to apply for security clearance (at the level of reliability status) with the National Research Council of Canada (NRC) in order to conduct fieldwork. At McGill, this may be of particular interest to radio astronomy observers and instrumentalists who are members of the CHIME or CHORD collaborations-both of which have telescope cores at the Dominion Radio Astrophysical Observatory (DRAO).

        It is possible to visit NRC sites without security clearance, provided such a visitor wears a guest badge and is accompanied at all times. This is not practical for deployments that may include imperfectly overlapping travel windows by various participants. Additionally, at the DRAO, a lack of security clearance makes it difficult to access the Blockhouse, a widely used lab outbuilding. For reasons like these, in collaborations like CHORD, it is strongly discouraged for collaborators without security clearance plan to work on site. 

        Once obtained, security clearance (at the level of reliability status) is valid for 10 years. The process is a bit bureaucratic but fairly straightforward for Canadian citizens who have not spent more than more than six cumulative months outside Canada in the five years up to the time of applying. Others should be prepared to enumerate the duration and reason for all travel outside Canada in the past five years. In the case of an applicant who has not known anyone in Canada for at least five years, the NRC is anecdotally unlikely to challenge the validity of character references the applicant has known for at least five years but who are from other countries.
        
        Non--Canadian citizens must also obtain a home country police report, which may require a second set of fingerprints and forwarding the report to the NRC via a physical mail courier. 

        It is prudent to start this application process several months before any planned travel to site. Fingerprints can be rejected due to illegibility, and, especially for home country police reports which require ink-based (as opposed to digital) fingerprints where an agent directs the applicant's hand, it can be difficult to directly control the success of the impression process. It may be necessary to return to the impression-enabled facility. As an example of how needing to repeat the fingerprinting process may be inconvenient, the default NRC fingerprinting location in Montr\'eal as of August 2025 was accessible via a 20-minute bus ride west of Station Namur on the Orange line. 

        The author of this section took every possible step to make the process go as quickly as possible: responding to all associated NRC emails and filling all forms within an hour of receiving them, making fingerprinting appointments on as short a timescale as possible, communicating to character references that (while understanding that they are very busy) it would be much appreciated if they could return the form on the sooner side because the process is unpredictable and long), and it still took about six weeks to receive clearance. 

        Once you have clearance, reach out to the project manager of your experiment, who will help you contact the appropriate people at the site you intend to work on to get the required NRC badge and pick a PIN for on-site buildings that require card access.

        Fun fact: once you have a NRC badge, since it is technically a government-issued photo ID, you can technically show it to board domestic flights in Canada. This can be handy if you prefer to minimize spurious commentary on any non-Canadian passport you might hold. One caveat is that, since your date of birth is not printed on the ID, some gate or security agents will not accept it as standalone ID. The author of this section has employed this strategy successfully at YUL and YLW but was unsuccessful using it at YYT.

       

